\section*{第 \ref{ch:b2:euler-bernoulli} 章——理想流体：欧拉与伯努利}

\begin{solution}{exo:b2:euler-bernoulli:1}
$v_2 = 1.5 \times 4 = \qty{6.0}{m/s}$，$P_2 = 3.0 \times 10^5 - 500(36 - 2.25) =
\qty{2.83}{bar}$。在 \qty{2}{cm} 处：$v = \qty{24}{m/s}$，$P = 3.0 \times 10^5 - 500
(576 - 2.25) = \qty{0.13}{bar}$——接近蒸气压。负值
说明所设流量不可能实现：水发生空化，
流动随之受阻。
\end{solution}

\begin{solution}{exo:b2:euler-bernoulli:2}
$v = \sqrt{2 \times 15000/0.74} = \qty{201}{m/s}$；指示空速 $\sqrt{2 \times 15000/
1.22} = \qty{157}{m/s}$。指示空速反映的是动压，
也就是机翼所感受到的量：失速速度是固定的\emph{指示}
速度，与高度无关。
\end{solution}

\begin{solution}{exo:b2:euler-bernoulli:3}
深度 \qty{2.0}{m}：$v = \sqrt{2g \times 2} = \qty{6.3}{m/s}$；$D_V = 0.6 \times 10^{-4}
\times 6.3 = \qty{0.38}{L/s}$；下落时间 $\sqrt{2/g} = \qty{0.45}{s}$，射程
\qty{2.8}{m}。射程 $= 2\sqrt{h(H - h)}$ 关于 $h \leftrightarrow H -
h$ 对称：深度 \qty{1.0}{m}（离地面 \qty{2}{m}）的小孔，其水流落在
同一点。
\end{solution}

\begin{solution}{exo:b2:euler-bernoulli:4}
$\Delta P = \tfrac12 \times 1.2 \times 900 = \qty{540}{Pa}$，外侧压强较低：
合力 \qty{54}{kN}，方向\emph{向上}；屋顶重 \qty{49}{kN}，屋顶会
被掀起。（实际屋顶往往先在边缘处破坏，那里的气流发生分离。）
\end{solution}

\begin{solution}{exo:b2:euler-bernoulli:5}
(a) $\tan\theta = 3/9.81$，$\theta = \qty{17}{\degree}$；$2 \times 0.306 = \qty{0.61}{m}$
（后部更高）。(b) $\Delta P = \rho aL = \qty{6.0}{kPa}$。(c) $\tan\theta =
0.61$，高度差 \qty{1.22}{m}：前端比平均液面高出
\qty{0.61}{m}，平均液面为 \qty{0.40}{m}：\qty{1.01}{m} $>$ \qty{0.8}{m}，液体会碰到盖子。(d) $a = v^2/R =
\qty{5.6}{m/s^2}$：$\theta = \qty{30}{\degree}$，外侧更高。
\end{solution}

\begin{solution}{exo:b2:euler-bernoulli:6}
(a) 旋转抛物面 $z = z_0 + \Omega^2r^2/2g$，$\Omega = 4\pi = \qty{12.6}{rad/s}$。(b)
$\Omega^2R^2/2g = 158 \times 0.0225/19.6 = \qty{0.18}{m}$。(c) 抛物面在圆盘上的平均高度
恰为中间值：中心 $0.20 - 0.09 = \qty{0.11}{m}$，边缘
\qty{0.29}{m}。(d) 中心恰好干涸的条件是 $\Omega^2R^2/4g = h_0$：$\Omega = \sqrt{4gh_0}/R =
\qty{18.7}{rad/s} = 3.0$ 转每秒。(e) 抛物面可将
平行光线会聚到一点，焦距为 $f = g/2\Omega^2$；液体镜面望远镜
让水银旋转而成。
\end{solution}

\begin{solution}{exo:b2:euler-bernoulli:7}
(a) $v = \sqrt{2g \times 3} = \qty{7.7}{m/s}$；$D_V = \pi \times 10^{-4} \times 7.7 =
\qty{2.4}{L/s}$。(b) $P = P_0 - \rho g(2) - \tfrac12\rho v^2 = 100 - 19.6 - 29.4 =
\qty{51}{kPa}$。(c) $P_0 - \rho gh_c - \tfrac12\rho v^2 \ge \qty{2.3}{kPa}$：$h_c \le
\qty{7.0}{m}$。(d) 空气可压缩且很轻：没有连续的
液柱来传递压强，两侧的水只是保持静止。
\end{solution}

\begin{solution}{exo:b2:euler-bernoulli:8}
(a) $\Delta P = mg/S = \qty{654}{Pa}$；$\tfrac12\rho v^2 = \qty{2.2}{kPa}$，$C_L = 0.30$。
(b) $v_{\text{上}}^2 = 58^2 + 2 \times 654/1.22$：\qty{66.6}{m/s}。(c) $F_L' =
\qty{981}{N/m} = \rho v\Gamma$：$\Gamma = \qty{13.4}{m^2/s}$。(d) $v = \sqrt{2mg/
\rho SC_L} = \sqrt{19620/25.6} = \qty{28}{m/s}$。
\end{solution}

\begin{solution}{exo:b2:euler-bernoulli:9}
(a) $v_2 = \qty{2.5}{m/s}$；$\Delta P = \tfrac12 \times 1060 \times (6.25 - 0.25) =
\qty{3.2}{kPa}$，$P_2 = \qty{9.8}{kPa}$。(b) 仍高于 \qty{1}{kPa}：否。(c) 在 
\qty{0.05}{cm^2}处，$v = \qty{5}{m/s}$，$\Delta P = \qty{13}{kPa}$：$P_2 \approx 0$，低于
组织压强——管壁塌陷，血流中止，
压强恢复，管壁重新打开：这便是一种颤振（医生听到的杂音）。
\end{solution}

\begin{solution}{exo:b2:euler-bernoulli:10}
(a) $c_PT_0 = c_PT + v^2/2$，其中 $c_P = \gamma R/(\gamma - 1)M_{\text{摩尔}}$，且
$c^2 = \gamma RT/M_{\text{摩尔}}$：$v^2/2c_PT = \tfrac12(\gamma - 1)M^2$。(b) $M =
0.85$：$T_0 = 220 \times 1.145 = \qty{252}{K}$；$M = 2$：$220 \times 1.8 = \qty{396}{K}$；
$M = 25$：$220 \times 126 = \qty{28000}{K}$——毫无意义：空气发生离解，
并发生电离，$c_P$ 不再是常数（真实的滞止温度
可达数千开尔文）。(c) $P_0/P = (1 + \tfrac{\gamma - 1}2M^2)^{\gamma/(\gamma
- 1)} \approx 1 + \tfrac{\gamma}2M^2 + \tfrac{\gamma}8M^4$，且 $\gamma PM^2 = \rho v^2$：
$P_0 - P = \tfrac12\rho v^2(1 + M^2/4)$。若将 $v$ 按 $\sqrt{2(P_0 - P)/\rho}$ 来读取，
则其被高估的比例为 $\sqrt{1 + M^2/4} - 1 \approx 9\%$（$M = 0.85$ 时）。
\end{solution}

\begin{solution}{exo:b2:euler-bernoulli:11}
(a) 从液面（$v \approx 0$，$P_0$，高度 $h$）到出口（$v$，$P_0$，
高度 $0$），其中 $\int\partial_tv\,\dd s = \ell\,\dd v/\dd t$。(b) $\dd v/(V^2 -
v^2) = \dd t/2\ell$，$V^2 = 2gh$：$\operatorname{artanh}(v/V) = Vt/2\ell$。(c) $V =
\qty{9.9}{m/s}$；$\tanh x = 0.9$ 时 $x = 1.47$：$t = 2\ell x/V = \qty{15}{s}$。(d)
$\dd v/\dd t \approx \qty{-20}{m/s^2}$：$\Delta P = \rho\ell \times 20 = \qty{9.9}{bar}$。
水的可压缩性很小：快速关闭会激起压力
波，$\Delta P = \rho c\Delta v \approx \qty{140}{bar}$。
\end{solution}

\begin{solution}{exo:b2:euler-bernoulli:12}
(a) 定常、理想、不可压缩且无旋：整个流场
只有一个常数。(b) 在液面上 $P = P_0$：$\tfrac12\rho v^2 + \rho gz =
0$，$z = -\Gamma^2/8\pi^2gr^2$。(c) 径向欧拉方程：$\partial_rP = \rho\Omega^2r$，
竖直方向为流体静力学：液面满足 $z(r) - z(a) = \Omega^2(r^2 -
a^2)/2g$。(d) 浴缸：$\Omega = \Gamma/2\pi a^2 = \qty{191}{rad/s}$；$z(a) = -\Omega^2a^2/
2g = \qty{-4.7}{cm}$，涡核还要再下降 \qty{4.7}{cm}：总深度为 \qty{9.3}{cm}。
龙卷风：外侧一份 $\Delta P = \tfrac12\rho v_{\max}^2$，内侧再同样一份：
$\rho v_{\max}^2 = 1.2 \times 6400 = \qty{7.7}{kPa}$，中心处比环境压强低。
\end{solution}

\begin{solution}{pb:b2:euler-bernoulli:1}
\textbf{1.} $\rho g \times 150 = \qty{14.7}{bar}$ 高于大气压；$\tfrac12\rho gH_d^2
\times \qty{1}{m} = \qty{1.1e8}{N}$。

\textbf{2.} $\rho gH = \qty{19.6}{bar}$ 表压，\qty{20.6}{bar} 绝对压强。

\textbf{3.} $v = 60/7.07 = \qty{8.5}{m/s}$; $\tfrac12\rho v^2 = \qty{0.36}{bar}$; $P =
P_0 + \rho gH - \tfrac12\rho v^2 = \qty{20.3}{bar}$ 绝对压强。

\textbf{4.} $v_j = \sqrt{2gH} = \qty{62.6}{m/s}$; $s = 60/62.6 = \qty{0.96}{m^2}$, $d =
\qty{1.1}{m}$.

\textbf{5.} $\tfrac12\rho D_Vv_j^2 = \tfrac12\rho D_V(2gH) = \rho gD_VH = \qty{118}{MW}$.

\textbf{6.} $m = \rho S\ell = \qty{2.8e6}{kg}$, $E_k = \tfrac12mv^2 = \qty{1.0e8}{J}$;
渡越时间 $\ell/v = \qty{47}{s}$。

\textbf{7.} $\rho gH = (P - P_0) + \tfrac12\rho v^2$：\qty{19.3}{bar} 的压强，
以及 \qty{0.36}{bar} 的动能；在喷管中压强降到
$P_0$，而 $\tfrac12\rho v^2$ 升到 $\rho gH$；紧接其后，$P = P_0$。

\textbf{8.} $\Delta h = \lambda(\ell/D)v^2/2g = 0.015 \times 133 \times 3.67 = \qty{7.3}{m}$,
$3.7\%$ 的功率。

\textbf{9.} $v = \qty{19.1}{m/s}$, $v^2/2g = \qty{18.6}{m}$, $\Delta h = 0.015 \times 200 \times
18.6 = \qty{56}{m}$：$28\%$ 损失——节省钢材的代价，
每秒都要以能量来支付。

\textbf{10.} $v_t = 60/\pi = \qty{19.1}{m/s}$; $\Delta P = 500(19.1^2 - 8.49^2) =
\qty{146}{kPa}$.

\textbf{11.} $\Delta P = (\rho_{\text{Hg}} - \rho)g\,\Delta h$: $\Delta h = 146000/(12600
\times 9.81) = \qty{1.18}{m}$.

\textbf{12.} $\Delta P \propto v^2 \propto D_V^2$：$D_V \propto \sqrt{\Delta h}$；流量减半
则降为四分之一，即 \qty{0.30}{m}。

\textbf{13.} 突出的测压孔迎着或扰动流动，读到的只是部分
动压（即皮托管），而非静压。

\textbf{14.} $P = P_0 - \rho g \times 6 - \tfrac12\rho v^2 = 100 - 59 - 36 = \qty{5}{kPa}$:
接近蒸气压时，溶解的空气逸出，水可能
沸腾（即空化），使水柱断开。降低峰顶，或装设
空气阀与真空泵，或减小流量。

\textbf{15.} $D_V/A = 60/2 \times 10^6 = \qty{3e-5}{m/s} = \qty{11}{cm/h}$.

\textbf{16.} $A\,\dd h/\dd t = -s\sqrt{2gh}$: $\sqrt h = \sqrt{h_0} - (s/A)\sqrt{g/2}\,t$.

\textbf{17.} $\Delta t = (A/s)\sqrt{2/g}(\sqrt{200} - \sqrt{190}) = 2.09 \times 10^6 \times
0.45 \times 0.36 = \qty{3.4e5}{s} = \qty{3.9}{days}$；若设计流量恒定，
$10A/D_V = \qty{3.3e5}{s}$：二者几乎相同，喷管已按
托里拆利在 $h \approx H$ 时的流量定尺寸。

\textbf{18.} $\ell\,\dd v/\dd t = gH - v^2/2$: $v = V\tanh(Vt/2\ell)$, $V = \qty{62.6}{m/s}$,
时间常数 $2\ell/V = \qty{13}{s}$。

\textbf{19.} 在 $v = V/2$ 处：$\dd v/\dd t = (gH - V^2/8)/\ell = \qty{3.7}{m/s^2}$；
$P = P_0 + \rho gH - \tfrac12\rho v^2 - \rho\ell\,\dd v/\dd t = 1.0 + 19.6 - 4.9 - 14.7
= \qty{1.0}{bar}$：水头的大部分用来加速水柱，故
管底压强远低于其稳态值。

\textbf{20.} $|\dd v/\dd t| = 8.49/10 = \qty{0.85}{m/s^2}$: $\Delta P = \rho\ell \times 0.85 =
\qty{3.4}{bar}$.

\textbf{21.} $\Delta P = 1000 \times 1400 \times 8.49 = \qty{119}{bar}$，是
静压的六倍。波的往返时间为 $2\ell/c = \qty{0.57}{s}$：凡快于此的
关闭称为“突然”关闭，获得完整的儒可夫斯基
超压；较慢的关闭所得超压较小。

\textbf{22.} $Av = A_s\,\dd z/\dd t$.

\textbf{23.} 沿隧洞应用欧拉方程，从水库（压强
$P_0 + \rho g\times$水深）至竖井底部（压强 $P_0 + \rho g
(\text{水深} + z)$）：$\rho L\,\dd v/\dd t = -\rho gz$；其中 $v = (A_s/A)\dot z$：
$\ddot z + (gA/LA_s)z = 0$; $T = 2\pi\sqrt{LA_s/gA} = 2\pi\sqrt{2000 \times 100/
68.7} = \qty{340}{s} \approx \qty{5.7}{min}$.

\textbf{24.} $\tfrac12\rho ALv^2 = \tfrac12\rho gA_sz_{\max}^2$，其中取 $v = 60/7 =
\qty{8.6}{m/s}$: $z_{\max} = v\sqrt{AL/gA_s} = \qty{32}{m}$.

\textbf{25.} 坝：流体静力学，$\rho gh$。涡轮机、阀门关闭时：
同样为 $\rho gH$。压力管道运行时：伯努利方程，$\tfrac12\rho v^2$ 取自
静水头。水锤：可压缩性，$\rho c\Delta v$。
涌浪：沿隧洞的非定常欧拉方程，其振荡周期为
$2\pi\sqrt{LA_s/gA}$.
\end{solution}
